\subsection{Benchmarking and Validation}\label{subsec:num_benchmarking}

This section discusses how to verify numerical calculations and compare against other methods.

\subsubsection{Convergence tests}

Vary integration tolerance and check convergence:

\begin{equation}
\left| C_\ell^{(\text{tol}_1)} - C_\ell^{(\text{tol}_2)} \right| < 0.01 \times C_\ell^{(\text{tol}_1)}
\end{equation}

If not satisfied, increase tolerance (reduce error).

\subsubsection{Limber limit validation}

For large $\ell$ (high multipoles), the Limber approximation should match full 3D calculation:

\begin{equation}
\left| \frac{C_\ell^{(\text{Limber})} - C_\ell^{(3D)}}{C_\ell^{(\text{Limber})}} \right| < 0.01 \quad \text{for } \ell > 100
\end{equation}

Deviations indicate numerical issues or incorrect kernel implementation.

\subsubsection{Comparison with literature benchmarks}

- BLAST paper provides reference $C_\ell$ values for specific cosmologies
- LSST DESC produces public forecasts (available online)
- N-body simulations (e.g., Illustris, EAGLE) provide ground-truth power spectra

\subsubsection{Unit tests}

\textbf{Example}: Test that spherical Bessel orthonormality integral equals expected value

\begin{equation}
\text{Result} = \int_0^\infty r^2 j_\ell(kr) j_\ell(k'r) dr
\end{equation}

Should equal $\frac{\pi}{2k^2} \delta(k-k')$ (or close for numerical $k, k'$).

\subsubsection{Profiling and timing}

Use `@time` or `@benchmark` (BenchmarkTools.jl) to identify bottlenecks:

\begin{itemize}
    \item Profile entire pipeline (Bessel → kernel → convolution)
    \item Typical timings: Bessel evaluation (60\%), integration (30\%), overhead (10\%)
    \item Goal: < 1 second per $C_\ell$ calculation on modern hardware
\end{itemize}

\remark{
Consistent timing and convergence with tighter tolerances builds confidence in results. Always compare against a known-good method before trusting new code.
}

See \cref{ch:numerical} for more on numerical strategies and \cref{subsec:num_performance} for optimization.
